\section{Introduction}
\label{sec:intro}

Beginning in July 1968, John B.~Calhoun followed a population of house mice in a closed enclosure that he called
\emph{Universe~25}, in which food, water and nest space were in excess, predators and epidemic disease were
excluded, and emigration was impossible \cite{Calhoun1973}. From four founding pairs the colony grew roughly
exponentially and stopped growing at about 2200 animals, well below the 3840 that the nest boxes could shelter. The
last young to survive weaning were born shortly after the peak, and the population then aged out without ever
recovering. Along the way Calhoun described a breakdown of social behaviour: withdrawn males gathered in pools at
the centre of the floor, nursing females attacked and abandoned their young, and there appeared the unwounded,
non-reproducing ``beautiful ones.'' Small groups moved to uncrowded enclosures during the decline did not reproduce,
even with partners raised without crowding (Marsden's work, reported in Ref.~\cite{Calhoun1973}; see also
Ref.~\cite{Marsden1972}). Calhoun read the collapse as a failure of social roles, not of resources: ``density per se was not the major factor'' \cite{Calhoun1973}.

The interpretation of the experiment has always been contested. Its popular message, ``density causes pathology,''
departed from Calhoun's own emphasis on social roles and design \cite{RamsdenAdams2009,Ramsden2009,AdamsRamsden2024},
and studies of human crowding found weak or no effects of density per se \cite{Galle1972,Freedman1975}. In earlier
confined rodent populations with abundant food, growth stopped through reduced litter survival mediated by social
interference \cite{Southwick1955,Lidicker1965}, and the maternal transmission of individual quality produces
delayed density dependence in models of population cycles \cite{GinzburgTaneyhill1994,InchaustiGinzburg1998}. For modelling, the decisive limitation is
that Universe~25 was a single realization; Calhoun himself called the study ``observation and reconstruction of a
process; it was historical'' \cite{Calhoun1973}. A model tuned to the published curve therefore risks fitting one
sample path of a strongly stochastic process \cite{GwizdallaDus2026}. The recent agent-based simulation of
Gwizdałła and Duś \cite{GwizdallaDus2026} calibrates the growth phase against
the reported counts and produces the decline through deviation factors that switch on when the \emph{total}
population exceeds critical values; other accounts invoke the loss of uncertainty about the social hierarchy in a
game-theoretic setting \cite{Mantzaris2025} or fit maximum-entropy unimodal curves without a mechanistic model
\cite{Roman2025}.

We ask instead which \emph{local} interaction rules, with unlimited resources and no rule that depends on the total
population, generate robustly across stochastic realizations the part of Calhoun's sequence that runs from rapid
growth to extinction: deceleration with social deterioration, cessation of reproduction near the peak, a decline
without recovery, and differentiated classes of deteriorated animals. The latency phase and the growth from four
founding pairs are part of that sequence, but we do not claim them (Sec.~\ref{sec:discussion}).

Our core model is a lattice model in which each mouse carries an age, a sex and a continuous social state
$s_i\in[0,1]$. The social state deteriorates under local crowding, scales conception and pup survival, and is partly
inherited; movement is biased towards acquaintances. We add four local rules, each with a neutral value that
switches it off and recovers the core model, and each motivated by an observation of the original report (Sec.~\ref{sec:model}): a
juvenile window of social plasticity (R1), social learning of competence (R2), crowd aversion of deteriorated
animals (AV), and refuges of limited capacity preferred by deteriorated animals (R7). Following pattern-oriented
modelling \cite{Grimm2005}, we translate Calhoun's report into seventeen falsifiable patterns, each with an
operational criterion evaluated run by run, and we measure the discriminating power of every essential criterion
against four null models: no crowding damage; full inheritance without social learning; the long-lived demography
without the new rules; and the calibrated candidate without its four rules (Sec.~\ref{sec:methods}). Seven
hypotheses were preregistered before the production runs, and every deviation from that plan is reported.

Our main result is negative, and it is controlled by the null models: \emph{the envelope is generic; the rules add
the classes}. Crowding damage and a long-lived demography alone already produce the envelope of the collapse---a
peak well below capacity with free space, the end of natality near the maximum, the social failure of late-born
cohorts when they are scored at the end of the plasticity window, and extinction. These patterns verify that the
model has a boom and its end; they are not evidence for the rules. What the rules add, and most null models lack,
is the coexistence of crowded and isolated deteriorated classes. The persistence of the isolated class follows from
the refuge design, and the failure of transplanted late-phase animals to found a colony from the plasticity window;
we report both as consistency checks, not findings. Second, the simulated dynamics is a boom of one or two generations, not Calhoun's
multi-generational cascade: about 70\% of the births are to the founding animals, and the net reproductive number
falls from 3.8 to 0.4 in a single generation. Third, the full set of patterns holds in 99.5\% of the runs of the
reference candidate, but this overstates its robustness: with a small background mortality or with spread initial
ages it falls to 55--63\%, only through a generic criterion on the timing of the end of natality, while the patterns
that the rules add pass in every run (Sec.~\ref{sec:results}). The minimal set of rules (R1, R2 and R7) is minimal
only at the evaluation thresholds used.

The ODD description, the full tables of patterns and null models, and the robustness analyses are given in the
Supplemental Material; the calibration, the mean-field analysis and the remaining results are in an extended version, available with the full ODD description, the code and the data in the repository \repourl.
